\section{A uniform potential bound for the Gauduchon ansatz}
\label{sec:stw-czero}

Let $X$ be a connected compact complex manifold of dimension $n\geq2$. Fix a Hermitian form $\gamma_0$ and a Gauduchon form $\gamma$ on $X$. For smooth real $u$, let $\Gamma_u$ be the $(n-1,n-1)$-form in the scalar ansatz of \cref{thm:stw-volume}; when $\Gamma_u>0$, let $\eta_u$ be its positive root, $\eta_u^{n-1}=\Gamma_u$. We study $\eta_u^n=\mathrm{e}^{F+b}\gamma^n$ for prescribed smooth real $F$ and an unknown real scalar $b$.

We can prove the zeroth-order estimate in this equation without invoking its existence theorem. The argument uses a fixed linear Green operator and a local supporting-plane construction. The latter is the elementary mechanism behind the real Alexandrov--Bakelman--Pucci estimate. The supporting-plane construction gives a determinant estimate adapted to the Gauduchon ansatz.

\begin{proposition}\label{prop:stw-czero-reconstructed}
Fix $\gamma_0$ and $\gamma$ as in \cref{thm:stw-volume}. If a smooth pair $(u,b)$ satisfies that equation, with $\Gamma_u>0$ and $\sup_Xu=0$, then
\[
 |b|+\|u\|_{L^\infty(X)}\leq C.
\]
The constant depends only on the fixed geometry and an upper bound for $\|F\|_\infty$. It does not depend on derivatives of $F$ or on a bounded mass hypothesis.
\end{proposition}
\begin{proof}
We first bound the scalar. At a maximum of $u$, its first derivatives vanish and $\ddc u\leq0$, so $\Gamma_u\leq\gamma_0^{n-1}$. At a minimum the reverse inequality holds. The determinant of the positive cofactor matrix is monotone under the order of positive matrices; the positive-root volume is its $1/(n-1)$ power. Evaluating the equation at these two points gives
\[
 \inf_X\left(\log\frac{\gamma_0^n}{\gamma^n}-F\right)
 \leq b\leq
 \sup_X\left(\log\frac{\gamma_0^n}{\gamma^n}-F\right).
\]
This proves the scalar bound.

Write
\[
 \mathcal D_\gamma u
 =\ddc u\wedge\gamma^{n-2}
 +\operatorname{Re}\left(\frac{\mathrm i}{2\pi}
                \partial u\wedge\bar\partial\gamma^{n-2}\right)
\]
and consider the fixed real linear elliptic operator
\[
 Lu=\frac{\gamma\wedge\mathcal D_\gamma u}{\gamma^n}.
\]
Its principal part is $\ddc u\wedge\gamma^{n-1}/\gamma^n$, its other term is first order, and it has no zeroth-order term. Positivity of $\Gamma_u$ implies $Lu\geq-C_0$ for a fixed constant $C_0$.

We recall explicitly the standard linear Green fact used here. A fixed smooth real uniformly elliptic operator with no zeroth-order term on a connected compact manifold has a smooth strictly positive adjoint probability $\nu$. With the sign for which $L$ is nonpositive at a maximum, its normalized Green kernel can be written
\[
 \begin{gathered}
 v(x)-\int_Xv\,\mathrm{d}\nu
       =-\int_X K(x,y)Lv(y)\,\mathrm{d}\nu(y),\\
 K\geq0,\qquad \sup_x\int_XK(x,y)\,\mathrm{d}\nu(y)<+\infty.
 \end{gathered}
\]
Here the kernel is smooth off the diagonal and locally integrable. To obtain this version, construct the inverse on functions of zero $\nu$-mean by the Fredholm alternative. The local fundamental solution has a positive leading singularity after the displayed minus sign, so compactness bounds its negative part. Adding a constant to $K$ makes it nonnegative and does not change the formula, since $\int_XLv\,\mathrm{d}\nu=0$. The usual local Green estimates give the uniform integral bound. These are standard linear elliptic Green estimates; no nonlinear equation estimate is being used. The positive adjoint density follows from the maximum-principle argument already used in the proof of \cref{thm:gauduchon}.

Apply the formula to $u$ and use $Lu\geq-C_0$. At a maximum point, where $u=0$, it gives
\[
 -\int_Xu\,\mathrm{d}\nu\leq C_1.
\]
Because $u\leq0$ and $\nu$ has a fixed smooth positive density, this is a uniform $L^1$ bound in any fixed smooth volume form.

Let $M=\min_Xu$ and choose holomorphic coordinates on a ball centered at a minimum point. Compactness permits a fixed coordinate radius and uniform bounds for the metric coefficients and their derivatives; rescale the ball to the unit ball $B\subset\R^{2n}$. Write $x$ for its real coordinates, $D$ and $D^2$ for the Euclidean gradient and Hessian, $\mathrm{d}x$ for Euclidean volume, and $I$ for an identity matrix of the indicated size. Fix a small number $\varepsilon>0$, depending only on those bounds, and set
\[
 v(x)=u(x)+\varepsilon|x|^2.
\]
For any real vector $p$ with $|p|<\varepsilon/2$, the function $v(x)-p\cdot x$ has an interior minimum: its boundary values are at least $M+\varepsilon/2$, whereas its value at the center is $M$. At such a point,
\[
 Dv=p,\qquad D^2v\geq0,\qquad
 v\leq M+\varepsilon/2.
\]
Let $A\subset B$ be the set of all points satisfying
$|Dv|\leq\varepsilon/2$, $D^2v\geq0$, and $v\leq M+\varepsilon/2$.
Since $u\geq M$, this set lies in $\{|x|\leq1/\sqrt2\}$. It is closed there, hence compact in $B$, and its gradient image contains the ball of radius $\varepsilon/2$. The area formula for the smooth gradient map therefore gives
\begin{equation}\label{eq:stw-contact-area}
 |B_{\varepsilon/2}|
       \leq\int_A\det D^2v\,\mathrm{d}x.
\end{equation}
The Hessian is nonnegative on $A$, so no absolute value is needed in the determinant.

We bound this determinant pointwise on $A$. There,
\[
 |Du|\leq\frac52\varepsilon,
 \qquad D^2u\geq-2\varepsilon I.
\]
The complex Hessian has the corresponding uniform lower bound. Wedge with the fixed positive $\gamma^{n-2}$ and bound the first-order term using the gradient estimate. Let $Q_u$ be the matrix of $\Gamma_u$ in a fixed local frame of $(n-1,0)$-forms. Since $\gamma_0^{n-1}$ is strictly positive, choosing $\varepsilon$ small gives
\[
 Q_u\geq cI
\]
with a uniform $c>0$. The equation and the already proved bounds for $F+b$ bound $\det Q_u$ from above, because the determinant of this cofactor matrix is, up to a fixed positive frame factor, the $(n-1)$st power of $\eta_u^n$. All eigenvalues of $Q_u$ are therefore bounded above as well: their product is bounded and each of the other $n-1$ eigenvalues is at least $c$.

In particular $\gamma\wedge\Gamma_u/\gamma^n$ is uniformly bounded on $A$. Subtracting its fixed background term and its bounded first-order term bounds
\[
 \frac{\ddc u\wedge\gamma^{n-1}}{\gamma^n}
\]
from above. This is a positive fixed multiple of the $\gamma$-trace of the complex Hessian of $u$. The complex Hessian of $v$ is nonnegative on $A$, since the real Hessian is nonnegative. Metric comparison thus bounds its Euclidean complex trace from above. The real trace of $D^2v$ is four times that complex trace. It follows that
\[
 0\leq\det D^2v
 \leq\left(\frac{\operatorname{tr}_{\R}D^2v}{2n}\right)^{2n}
 \leq C_2\quad\hbox{on }A.
\]
Equation \eqref{eq:stw-contact-area} now gives $|A|\geq c_2>0$, where $c_2$ is independent of $u$. On $A$ we also have $u\leq v\leq M+\varepsilon/2$. Consequently the uniform $L^1$ bound gives
\[
 C_3\geq\int_A(-u)\,\mathrm{d}x
       \geq c_2(-M-\varepsilon/2).
\]
This bounds $-M$, and completes the proof.
\end{proof}

The estimate supplies the uniform potential bound needed in a continuity method for the Gauduchon ansatz. The second-order estimate and the rescaling argument below will supply the derivative bounds, retaining the first-order torsion terms.
